\documentclass[11pt]{article}
\usepackage[margin=1in]{geometry}
\usepackage[T1]{fontenc}
\usepackage[utf8]{inputenc}
\usepackage{lmodern}
\usepackage{enumitem}
\usepackage{hyperref}
\usepackage{titlesec}
\setlist[itemize]{leftmargin=1.5em}
\titleformat{\section}{\large\bfseries}{\thesection.}{0.5em}{}
\titleformat{\subsection}{\normalsize\bfseries}{\thesubsection.}{0.5em}{}

\title{Methods Outline and Writing Brief}
\author{}
\date{}

\begin{document}
\maketitle

\section{Study Design}
\textbf{Purpose:} State the paper's objective and the full analytical workflow.

\textbf{Key sentences:}
\begin{itemize}
\item We evaluated how hierarchical simplification of multichannel river networks alters network geometry and event-scale routing response.
\item The workflow combined real-case selection, mask-based network extraction, hierarchy detection, mask-based collapse and regeneration, hydraulic attribute reconstruction, equivalent-resistance analysis, and comparative routing simulations.
\end{itemize}

\textbf{Include:}
\begin{itemize}
\item real cases and synthetic support experiments
\item geometry and topology metrics
\item equivalent resistance
\item event-scale routing comparison
\end{itemize}

\section{Case Selection}
\textbf{Purpose:} Explain why the examples were chosen and how they were constrained.

\textbf{Key sentences:}
\begin{itemize}
\item Study cases were defined as multichannel river sections bounded by upstream and downstream gauging stations.
\item Cases were retained only where geomorphic complexity and hydrologic data availability supported event-scale comparative analysis.
\end{itemize}

\textbf{Include:}
\begin{itemize}
\item gauge-pair logic
\item SWORD matching
\item subdaily discharge availability
\item event suitability and quality screening
\item final retained study set
\end{itemize}

\textbf{Figure suggestion:} map of selected cases.

\section{Synthetic Support Experiments}
\textbf{Purpose:} Present synthetic networks as a methodological support component.

\textbf{Key sentences:}
\begin{itemize}
\item Synthetic networks were used as controlled methodological support to isolate the effects of topology, width partitioning, slope contrasts, sinuosity contrasts, and forcing characteristics.
\item These experiments supported interpretation of the real-network simplification results rather than replacing them.
\end{itemize}

\textbf{Include:}
\begin{itemize}
\item base corridor setup
\item looped and cross-linked perturbations
\item width-conserving branch partitioning
\item sensitivity dimensions
\end{itemize}

\textbf{Figure suggestion:} schematic of the synthetic setup.

\section{Mask Preparation and RivGraph Extraction}
\textbf{Purpose:} Describe the image-to-network conversion.

\textbf{Key sentences:}
\begin{itemize}
\item Source water masks were converted to binary rasters and manually cleaned where necessary before centerline and network extraction with RivGraph.
\item This produced centerline, node, and link geometries for each case.
\end{itemize}

\textbf{Include:}
\begin{itemize}
\item binary raster preparation
\item manual mask cleaning
\item RivGraph extraction
\end{itemize}

\section{Directed Graph Preparation}
\textbf{Purpose:} Explain how direction assignment was handled and frame manual review as optional.

\textbf{Key sentences:}
\begin{itemize}
\item Flow direction was assigned primarily from RivGraph outputs.
\item Where automated direction assignment was uncertain, manual review of link direction and inlet--outlet roles could be used as an additional quality-control step.
\item All downstream analyses required a valid directed network with one source and one sink.
\end{itemize}

\textbf{Include:}
\begin{itemize}
\item RivGraph flow assignment
\item optional manual review and correction
\item graph validation criteria
\end{itemize}

\section{Hierarchical Unit Detection}
\textbf{Purpose:} Define the core analysis object.

\textbf{Key sentences:}
\begin{itemize}
\item The network was decomposed into bifurcation--confluence units, each defined by one or more directed paths between an upstream bifurcation and a downstream confluence.
\item Units could be simple, multi-thread, or compound and nested.
\end{itemize}

\textbf{Include:}
\begin{itemize}
\item path-based unit definition
\item nested units
\item compound bubbles as broader multichannel zones
\end{itemize}

\textbf{Figure suggestion:} annotated unit example on a real network.

\section{Metrics and Collapse Ranking}
\textbf{Purpose:} Justify how candidate collapses were prioritized.

\textbf{Key sentences:}
\begin{itemize}
\item Candidate units were characterized using geometric persistence, path-width disparity, and internal topologic complexity.
\item Collapse ranking was based on path disparity, elongation, equivalent length, and topologic complexity.
\end{itemize}

\textbf{Include:}
\begin{itemize}
\item \texttt{path\_disparity\_width}
\item \texttt{equivalent\_length}
\item \texttt{elongation}
\item \texttt{topologic\_complexity\_score}
\item grouped candidates selected from ranked contiguous partitions
\end{itemize}

\section{Collapse and Variant Regeneration}
\textbf{Purpose:} Explain that simplification was implemented in mask space and regenerated with RivGraph.

\textbf{Key sentences:}
\begin{itemize}
\item Simplification was implemented by modifying the parent mask within selected collapse domains and rerunning RivGraph, rather than deleting graph elements directly.
\item This ensured that each simplified state was a regenerated network geometry.
\end{itemize}

\textbf{Include:}
\begin{itemize}
\item footprint construction
\item filling enclosed dry holes
\item independent and sequential experiments
\item realized versus anchored interpretations if used in the paper
\end{itemize}

\textbf{Figure suggestion:} original versus collapsed example.

\section{Width Reconstruction}
\textbf{Purpose:} Describe the width methodology clearly and concisely.

\textbf{Key sentences:}
\begin{itemize}
\item Widths were recomputed on regenerated variants using a two-bank distance-transform-based method.
\item For each regenerated link, total width was measured on the collapsed mask, wet width on the original wet mask, and dry width as the difference between the two.
\end{itemize}

\textbf{Include:}
\begin{itemize}
\item sample-based bank search
\item total, wet, and dry width families
\item why wet versus dry width matters after collapse
\end{itemize}

\textbf{Figure suggestion:} width-method figure, ideally combined with original and collapsed geometry.

\section{Hydraulic Context Reconstruction}
\textbf{Purpose:} Explain how slope and profile context were reconstructed.

\textbf{Key sentences:}
\begin{itemize}
\item Node water-surface-elevation context from matched external reference nodes was converted to a monotone or interpolated longitudinal profile for slope estimation.
\item Observed node elevation context and derived routing profile were treated as distinct quantities.
\end{itemize}

\textbf{Include:}
\begin{itemize}
\item SWORD node matching
\item profile-based slope derivation
\end{itemize}

\section{Geometry-Informed Equivalent Resistance}
\textbf{Purpose:} Position equivalent resistance as a core analytical layer.

\textbf{Key sentences:}
\begin{itemize}
\item Equivalent resistance was used as a geometry-informed predictor of the hydraulic effect of multichannel structures and their simplified replacements.
\item Edge-level resistances were computed from link geometry and reduced to equivalent corridor resistance under normalized head boundary conditions.
\item The principal experiment-level comparison metric was the ratio between realized child-corridor resistance and matched base-corridor resistance.
\end{itemize}

\textbf{Include:}
\begin{itemize}
\item state-level bubble resistance
\item experiment-level matched comparison
\item interpretation of \texttt{Lambda\_matched}
\end{itemize}

\textbf{Figure suggestion:} schematic of bubble reduction and matched corridor comparison.

\section{Event Forcing and Routing Comparison}
\textbf{Purpose:} Explain the comparative hydrodynamic experiment.

\textbf{Key sentences:}
\begin{itemize}
\item All states from a given example were forced with the same observed event hydrograph and routed with the same parameterization.
\item Results were interpreted as relative differences in timing, attenuation, and hydrograph shape across variants.
\end{itemize}

\textbf{Include:}
\begin{itemize}
\item shared forcing per example
\item shared routing parameterization
\item comparative rather than state-specific calibration framing
\item peak timing, attenuation, rise time, fall time, and hydrograph shape
\end{itemize}

\section{Suggested Figure Set}
\begin{itemize}
\item Figure 1: map of selected real cases
\item Figure 2: synthetic network setup and sensitivity dimensions
\item Figure 3: hierarchy definition on a real reviewed network
\item Figure 4: collapse and regeneration workflow on one example
\item Figure 5: width methodology on original versus collapsed geometry
\item Figure 6: equivalent-resistance schematic and matched comparison concept
\end{itemize}

\end{document}
